% Tabla de contenido con una cabecera por capítulo y páginas automáticas.
\addto\captionsspanish{\renewcommand{\contentsname}{TABLA DE CONTENIDO}}

\makeatletter
\newcommand{\l@chaptergroup}[2]{%
  \addvspace{1em}%
  \@dottedtocline{0}{0em}{0em}{\bfseries #1}{\bfseries #2}}
\newcommand{\toclevel@chaptergroup}{0}
\renewcommand{\l@chapter}[2]{%
  \@dottedtocline{0}{0em}{1.8em}{\bfseries #1}{\bfseries #2}}

% Conserva la numeración normal de book y sus enlaces; solo añade la
% cabecera CAPÍTULO N. antes del título numerado en la tabla de contenido.
\patchcmd{\@chapter}
  {\addcontentsline{toc}{chapter}{\protect\numberline{\thechapter}#1}}
  {\addcontentsline{toc}{chaptergroup}{\protect\MakeUppercase{\@chapapp}\ \thechapter.}%
   \addcontentsline{toc}{chapter}{\protect\numberline{\thechapter}#1}}
  {}
  {\PackageError{tesis}{No se pudo configurar la tabla de contenido}%
    {Comprueba la definicion de chapter en la clase book.}}
\makeatother
